\section{Dimensional Compression Ladders}
\label{sec:compression_results}

Dimensional compression evaluates representation retention when high-dimensional embeddings ($d \in \{768, 1024, 1152\}$) are compressed into compact subspaces $k \in \{512, 256, 128, 64\}$. We evaluate five compression algorithms: coordinate prefix slicing (\texttt{prefix}), unweighted principal component projection (\texttt{pca}), truncated variance-equalizing whitening (\texttt{whitening}), adaptive spectral tempering (\texttt{spectemp}), and uniform random coordinate subsampling (\texttt{random\_truncation}).

Tables~\ref{tab:arguana_compression_ndcg} to~\ref{tab:stsb_compression_pearson} in Appendix~\ref{app:app2} report absolute scores and retention for every model, method, and $k$. Retention is the compressed score divided by the same model's full-dimension baseline score, so values above 100\% mean the compressed embedding scored higher than the raw one.

\paragraph{Retrieval loses far more than similarity under compression.}
For dedicated encoders, the loss at small $k$ depends strongly on the task. On FiQA2018 and SCIDOCS, even the best method keeps only about three quarters of nDCG@10 at $k=64$, and about 90\% on ArguAna. On all three STS datasets, every method keeps at least 94\% of Spearman $\rho$ at $k=64$.

\paragraph{Whitening is costly at large $k$ and recovers at small $k$.}
For dedicated encoders at $k=512$, whitening is the weakest method on every dataset except STS22.v2 Spearman $\rho$, while prefix slicing, PCA, and SpecTemp stay close to the native score. The gap closes as $k$ shrinks. At $k=64$, whitening is level with PCA on FiQA2018 and SCIDOCS and ahead of PCA on STS22.v2 and STSBenchmark.

\paragraph{Prefix slicing holds up on similarity but not on retrieval.}
Prefix slicing keeps nearly all of the dedicated score at $k=512$ on every dataset. At $k=64$ it is the weakest method apart from random truncation on all three IR datasets, while on STS it still keeps at least 96\% of Spearman $\rho$. On retrieval, EmbeddingGemma's prefix loses much more than Qwen3-Embedding's at $k=64$ on every dataset.

\paragraph{PCA and SpecTemp trade places at small $k$.}
For dedicated encoders, PCA and SpecTemp both track the native score closely at $k=512$, and on SCIDOCS both slightly exceed it for EmbeddingGemma. At $k=64$, PCA keeps more on ArguAna, while SpecTemp keeps more on FiQA2018 and SCIDOCS.

\paragraph{Random truncation is weak on retrieval but not on similarity.}
For dedicated encoders, random coordinate subsets keep the least nDCG@10 of all methods at $k=64$ on every IR dataset. On STS they keep at least 95\% of Spearman $\rho$, and on STS22.v2 they keep more than any other method.

\paragraph{The ranking of methods flips for base models.}
For base models, whitening at $k=512$ scores above the native nDCG@10 or Spearman $\rho$ for every model and pooling mode on every dataset, and above SpecTemp in every case. On retrieval this advantage shrinks with $k$, and the base average falls to or below the native score by $k=64$. On STS, the base average for whitening stays above the native Spearman $\rho$ down to $k=64$. Very large retention values occur where the native score is close to zero: Qwen3.5-0.8B-Base with last-token pooling has a native FiQA2018 nDCG@10 of 0.0026, so whitening's 0.0060 at $k=512$ appears as 233\% retention.

\paragraph{STS22.v2 behaves differently at $k=512$.}
On STS22.v2 at $k=512$, whitening keeps on average only about half of the dedicated Pearson $r$, while its Spearman $\rho$ stays near the native score. At $k=256$ and $k=128$, the dedicated average for whitening is at or above the native Pearson $r$. Also at $k=512$, SpecTemp gives exactly the same scores as PCA for every model, which matches the zero tempering exponent reported for EmbeddingGemma (Table~\ref{tab:spectemp_dynamics}).
